% 2010 年天津卷文科第 12 题：正视图、侧视图与俯视图。
% 来源 PDF：普通高考/2010/2010天津文.pdf（上游勘误版），SHA-256 e6b06813e17b7c332b651b8784c8fab0fcad0c8666ea6032a2958917c55307c9。
% 原页：第 1 页（A3 四栏）第 3 栏上部，书页码 637。
% 重绘依据：正视图为 2×1 矩形、中分；侧视图为 1×1 矩形；俯视图为梯形，上底 2、下底 1、高 2。
\begin{tikzpicture}[semithick, every node/.style={inner sep=1.2pt, font=\small}]
  % 正视图
  \draw (0,4.2) rectangle (2,5.2);
  \draw (1,4.2) -- (1,5.2);
  \node[below] at (1,3.7) {正视图};
  \draw[<->] (0,5.45) -- (1,5.45) node[midway,above] {$1$};
  \draw[<->] (1,5.45) -- (2,5.45) node[midway,above] {$1$};
  \draw[<->] (-0.4,4.2) -- (-0.4,5.2) node[midway,left] {$1$};
  % 侧视图
  \draw (3.4,4.2) rectangle (4.4,5.2);
  \node[below] at (3.9,3.7) {侧视图};
  \draw[<->] (3.4,5.45) -- (4.4,5.45) node[midway,above] {$1$};
  % 俯视图
  \draw (0,0) -- (1,0) -- (2,2) -- (0,2) -- cycle;
  \draw[<->] (0,-0.35) -- (1,-0.35) node[midway,below] {$1$};
  \draw[<->] (-0.4,0) -- (-0.4,2) node[midway,left] {$2$};
  \node[below] at (1.1,-0.95) {俯视图};
\end{tikzpicture}
